\referenceschapter

\begin{thebibliography}{9}
\bibitem{deBerg2008} Mark de Berg, Otfried Cheong, Marc van Kreveld and Mark Overmars,
  \emph{Computational Geometry: Algorithms and Applications}, 3rd edition,
  Springer, 2008. The primary source for the algorithms formalised here,
  including planar arrangements, the sweep-line method, point location and
  convex hulls.

\bibitem{cgal} The \link{https://www.cgal.org/}{CGAL} project, \emph{The Computational
  Geometry Algorithms Library}. Influential on the design of this library,
  particularly its separation of geometric algorithms from the arithmetic
  used to instantiate them, and its treatment of exact predicates.

\bibitem{kettner2008} Lutz Kettner, Kurt Mehlhorn, Sylvain Pion, Stefan Schirra and Chee Yap,
  \emph{Classroom examples of robustness problems in geometric computations},
  Computational Geometry \textbf{40}(1):61--78, 2008.
  \link{https://doi.org/10.1016/j.comgeo.2007.06.003}{doi:10.1016/j.comgeo.2007.06.003}.
  Exhibits concrete inputs on which floating-point implementations of convex
  hull and Delaunay triangulation fail, and analyses the geometry of the
  floating-point orientation predicate. This is the empirical case for
  Chapter~\ref{chap:numbers}.

\bibitem{flocq} Sylvie Boldo and Guillaume Melquiond, \emph{Flocq: A Unified Library for
  Proving Floating-Point Algorithms in Coq}, ARITH 20, pp.\ 243--252, IEEE, 2011.
  \link{https://doi.org/10.1109/ARITH.2011.40}{doi:10.1109/ARITH.2011.40}.
  A formalization of IEEE 754 floating-point arithmetic in Coq, and a model for
  how hardware floats could back the fast approximations of
  Chapter~\ref{chap:numbers}.

\bibitem{mathlib2020} The mathlib Community, \emph{The Lean Mathematical Library}, CPP 2020,
  pp.\ 367--381.
  \link{https://doi.org/10.1145/3372885.3373824}{doi:10.1145/3372885.3373824}.
  The source of the algebraic and order hierarchies used throughout, and of the
  \code{Set}-level definitions against which the geometric results are
  ultimately stated.

\bibitem{lean4} Leonardo de Moura and Sebastian Ullrich, \emph{The Lean 4 Theorem Prover
  and Programming Language}, CADE 28, LNCS 12699, pp.\ 625--635, Springer, 2021.
  \link{https://doi.org/10.1007/978-3-030-79876-5_37}{doi:10.1007/978-3-030-79876-5\_37}.
  Describes the system, including its treatment of Lean as a programming
  language, which is what makes it reasonable to expect these algorithms to run
  as well as to be proved correct.

\bibitem{rau2020} Martin Rau and Tobias Nipkow, \emph{Verification of Closest Pair of Points
  Algorithms}, IJCAR 2020, LNCS 12167, pp.\ 341--357.
  \link{https://doi.org/10.1007/978-3-030-51054-1_20}{doi:10.1007/978-3-030-51054-1\_20}.
  The closest published analogue to this project: divide-and-conquer geometric
  algorithms verified for both functional correctness and asymptotic running
  time, with executable code extracted and benchmarked against handwritten
  implementations. A useful reference point for the choices discussed in
  Chapter~\ref{chap:arrangements}.

\bibitem{nipkow2020} Tobias Nipkow, Manuel Eberl and Maximilian P.\ L.\ Haslbeck,
  \emph{Verified Textbook Algorithms: A Biased Survey}, ATVA 2020, LNCS 12302.
  \link{https://doi.org/10.1007/978-3-030-59152-6_2}{doi:10.1007/978-3-030-59152-6\_2}.
  Surveys what has been achieved in verifying standard algorithms, covering both
  correctness and running time, and is the most direct discussion of the
  trade-offs weighed in Chapter~\ref{chap:arrangements}: how much of a
  specification to state abstractly, when to refine a clean definition into an
  efficient implementation, and what is lost either way.

\bibitem{knuth1992} Donald E.\ Knuth, \emph{Axioms and Hulls}, Lecture Notes in Computer
  Science 606, Springer, 1992.
  \link{https://doi.org/10.1007/3-540-55611-7}{doi:10.1007/3-540-55611-7}.
  The source of the counterclockwise axioms in
  Chapter~\ref{chap:primitives}. Knuth replaces coordinates by a ternary
  predicate satisfying five axioms, and develops convex hull algorithms against
  the axioms alone, so that any predicate satisfying them --- whatever
  arithmetic computes it --- inherits the algorithms' correctness.

\bibitem{chlipala2013} Adam Chlipala, \emph{Certified Programming with Dependent Types: A
  Pragmatic Introduction to the Coq Proof Assistant}, MIT Press, 2013.
  \link{https://adam.chlipala.net/cpdt/}{adam.chlipala.net/cpdt}.
  Written for Coq rather than Lean, but its subject is proof engineering rather
  than any particular prover: how to structure a development so that it survives
  change. Its treatment of dependently typed programming --- when to push
  correctness into the type of a function and when to keep it in a separate
  theorem --- is the question Chapter~\ref{chap:arrangements} has to answer for
  each algorithm.
\end{thebibliography}

% vim: ts=2 sw=2 et ai
